% from ega0-21-fr.tex
\oldpage[0\textsubscript{IV-1}]{163}
\begin{corollary}[21.4.4]
\label{0.21.4.4-fr}
Soient $k$ un corps de caractéristique $p>0$, $K$ une extension de $k$.
Pour que $\Omega^1_{K/k}=0$, il faut et il suffit que $K=k(K^p)$.
En particulier, pour que $\Omega^1_K=0$, il faut et il suffit que $K$ soit
un corps parfait.
\end{corollary}

\begin{proof}
En effet, si $(x_\alpha)$ est une $p$-base de $K$ sur $k$, les
$d_{K/k}(x_\alpha)$ forment une base de l'espace vectoriel $\Omega^1_{K/k}$
(21.2.5).
\end{proof}

\begin{theorem}[21.4.5]
\label{0.21.4.5-fr}
Soient $k$ un corps de caractéristique $p>0$, $K$ une extension de $k$,
$(x_\alpha)$ une famille d'éléments de $K$. Pour que $(x_\alpha)$ soit
$p$-libre sur $k$ (resp. une famille de $p$-générateurs de $K$ sur $k$,
resp. une $p$-base de $K$ sur $k$), il faut et il suffit que la famille
$(d_{K/k}(x_\alpha))$ soit une famille libre (resp. un système de
générateurs, resp. une base) dans l'espace vectoriel $\Omega^1_{K/k}$.
\end{theorem}

\begin{proof}
